% Generated draft: requires mathematical and source review.
\documentclass[11pt]{article}
\usepackage[margin=25mm]{geometry}
\usepackage{fontspec}
\setmainfont{Latin Modern Roman}
\usepackage{amsmath,amssymb,mathtools,unicode-math}
\setmathfont{Latin Modern Math}
\usepackage{enumitem,needspace,xurl,hyperref}
\setlength{\parindent}{0pt}
\setlength{\parskip}{5pt}
\setlength{\emergencystretch}{3em}
\newcommand{\N}{\mathbb{N}}
\newcommand{\Z}{\mathbb{Z}}
\newcommand{\Q}{\mathbb{Q}}
\newcommand{\R}{\mathbb{R}}
\newcommand{\C}{\mathbb{C}}
\newcommand{\F}{\mathbb{F}}
\newcommand{\A}{\mathbb{A}}
\newcommand{\PP}{\mathbb{P}}
\newcommand{\E}{\mathbb{E}}
\newcommand{\eps}{\varepsilon}
\newcommand{\dd}{\,\mathrm{d}}
\begin{document}
\section*{UnsolvedMath ID: 800598}
\textbf{Source problem code:} catalogue entry
\subsection{Definitions and mathematical statement}
Let \(X=\{0,1\}^{\N}\), equipped with the product topology in which \(\{0,1\}\) is discrete. Write \([\N]^{\aleph_0}\) for the set of infinite subsets of \(\N\). For \(a,b\subseteq\N\), write \(a\subseteq^*b\) if \(a\setminus b\) is finite. A family \(Y\subseteq[\N]^{\aleph_0}\) is \emph{linearly refinable} if there is a choice of infinite sets \(\widehat y\subseteq y\), one for each \(y\in Y\), such that
\[
\forall y,z\in Y\quad
\bigl(\widehat y\subseteq^*\widehat z\ \lor\ \widehat z\subseteq^*\widehat y\bigr).
\]
A countable open cover \(\mathcal U=\{U_n:n\in\N\}\) of \(X\) is \emph{large} if every \(x\in X\) belongs to infinitely many members of \(\mathcal U\). Define its incidence sets by
\[
a_x=\{n\in\N:x\in U_n\}\qquad(x\in X).
\]
The cover is a \(\tau^*\)-cover if it is large and \(\{a_x:x\in X\}\) is linearly refinable. Let \(\mathrm T^*(X)\) denote the collection of countable open \(\tau^*\)-covers of \(X\). An infinite open cover \(\mathcal V\) is a \(\gamma\)-cover if
\[
\forall x\in X\quad \{V\in\mathcal V:x\notin V\}\text{ is finite}.
\]
Let \(\Gamma(X)\) denote the collection of countable open \(\gamma\)-covers. The notation \(\binom{\mathrm T^*}{\Gamma}\) means that every member of \(\mathrm T^*(X)\) contains a subfamily belonging to \(\Gamma(X)\).

\emph{Question (source Problem 26).} Does the Cantor space \(X\) satisfy
\[
\forall\mathcal U\in\mathrm T^*(X)\quad
\exists\mathcal V\subseteq\mathcal U\quad
\mathcal V\in\Gamma(X)?
\]
Equivalently, for a cover enumerated without repetitions, must there exist an infinite \(I\subseteq\N\) such that
\[
\forall x\in X\quad I\subseteq^*a_x?
\]
Here the infinite refinements witnessing linear refinability may depend on \(x\), whereas the desired set \(I\) must work simultaneously for every \(x\in X\).

\emph{Unresolved conventions and context.} The definitions of large cover, \(\Gamma\), and the binomial selection notation above use their standard meanings; the supplied record does not spell them out. This formulation treats covers as families of distinct sets and does not impose the additional condition \(X\notin\mathcal U\). The intended convention about proper cover members and repeated terms requires confirmation from the source. As background, the record also asserts that every analytic space satisfies \(\binom{\mathrm T}{\Gamma}\), with reference label \([\mathrm{tau}]\). It does not define \(\mathrm T\), specify the scope of that assertion, or supply the bibliography entry for that label; the assertion is retained as supplied context, not independently verified.
\subsection{Short English statement}
Does every countable open \(\tau^*\)-cover of the Cantor space contain an infinite subcover in which each point belongs to all but finitely many members? A \(\tau^*\)-cover allows each point's infinite incidence set to be thinned so that the resulting sets are linearly ordered modulo finite differences. The registry labels this problem open; cover conventions and unresolved reference labels require human review.
\subsection{Sources}
\begin{itemize}
\item Local UnsolvedMath registry, permanent ID \(800598\), exact source problem code \texttt{catalogue entry}, AMR Open Problem Lists. The record attributes the problem to Tsaban, with proposed year 2006, and identifies it as Problem 26, with attribution/context label \texttt{[tautau]}. Its supplied status is \emph{open}; the record was updated on 2026-10-16. This is supplied registry metadata, not an independently verified status assessment.
\item Tsaban, \emph{Selection Principles and Special Sets of Reals: Open Problems} (2006), Problem 26: \url{https://arxiv.org/abs/math/0609601}. The record reports access on 2026-10-16 and extraction by \texttt{arxiv\_latex}. These are supplied provenance details; no independent source visit or verification is claimed.
\item The supplied contextual assertion about analytic spaces cites \texttt{[tau]}, and the problem title/context uses \texttt{[tautau]}. Corresponding bibliography entries are absent from the supplied record, so their authorship and publication details remain unresolved.
\end{itemize}
\end{document}
